\chapter{Resultados y discusión}\label{cap:resultados}

\section{Estado funcional alcanzado}

\begin{itemize}[leftmargin=*,itemsep=0.2em]
  \item Backend FastAPI con $\geq 5$ endpoints productivos (predict,
        verify, save, delete, billing).
  \item Frontend React con autenticación, dashboard, historial,
        gestión de claves API, suscripciones Stripe, panel
        administrativo, multi-idioma, cierre por inactividad.
  \item Extensión MV3 con login, análisis y verificación.
  \item Supabase con migraciones SQL versionadas y RLS activa.
      \item Clasificador NLI FEVER exportado a ONNX e integrado en el
                        backend mediante \texttt{FEVER\_MODEL\_PATH}.
\end{itemize}

\section{Métricas técnicas}

\begin{table}[h]
  \centering
  \begin{tabular}{lcc}
    \toprule
    Componente & Líneas de código & Tests \\
    \midrule
      Backend FastAPI         & 2\,696 & 38 \\
      Frontend React          & 8\,549 & build y tests autom. \\
      Extensión navegador     & 2\,223 & validación manual \\
      Módulo experimental FEVER & 869 & evaluación estadística \\
    \bottomrule
  \end{tabular}
  \caption{Volumen y cobertura de pruebas por componente.}
  \label{tab:metricas}
\end{table}

\section{Latencias observadas}

La validación local previa a la entrega confirmó el arranque del
backend y del frontend de desarrollo. El \emph{healthcheck} del backend
respondió correctamente en entorno local; los endpoints que consumen
cuota y guardan historial requieren una sesión Supabase válida, por lo
que se validan mediante tests de backend y pruebas manuales con usuario
autenticado. No se ha realizado un benchmark p50/p95 exhaustivo, de
modo que las latencias deben interpretarse como comprobación funcional,
no como caracterización de rendimiento en producción.

\section{Calidad del agente de verificación}

El clasificador integrado en el agente es DeBERTa-v3-small entrenado
sobre FEVER-NLI y exportado a ONNX. La tabla~\ref{tab:fever-integrado}
resume la comparación frente al baseline TF-IDF + regresión logística.

\begin{table}[h]
      \centering
      \begin{tabular}{lcccc}
            \toprule
            Modelo & Acc. & F1 macro & NLL & ECE \\
            \midrule
            TF-IDF + LR & 0.5081 & 0.4842 & 1.0897 & 0.1585 \\
            DeBERTa-v3-small ONNX & 0.7281 & 0.7231 & 0.8070 & 0.1277 \\
            \bottomrule
      \end{tabular}
      \caption{Resultados FEVER-NLI usados para seleccionar el clasificador integrado.}
      \label{tab:fever-integrado}
\end{table}

Además de mejorar la calidad estadística, el formato ONNX permite que
la inferencia del backend no dependa de cargar \texttt{torch}. Si el
artefacto \texttt{model.onnx} no está disponible, el backend conserva
un clasificador determinista de respaldo para que el entorno de pruebas
arranque de forma reproducible.

\section{Limitaciones}

\begin{itemize}[leftmargin=*,itemsep=0.2em]
  \item No se ha implementado escalado horizontal automático.
  \item No se ha integrado un sistema de \emph{tracing} distribuido
        (OpenTelemetry).
  \item La extensión está probada en navegadores Chromium; falta
        validación exhaustiva en Firefox.
  \item Las probabilidades del clasificador NLI no están calibradas
        \emph{post-hoc}; el ECE mejora al baseline, pero se interpreta
        la confianza como señal relativa.
\end{itemize}

\section{Aprendizajes}

\begin{itemize}[leftmargin=*,itemsep=0.2em]
  \item La separación estricta UI / Estado / Servicios ha facilitado
        las pruebas y la incorporación posterior del agente sin
        introducir regresiones.
  \item RLS en Supabase reduce significativamente el riesgo de fugas
        accidentales de datos entre usuarios, pero exige disciplina
        en migraciones y revisiones de código.
  \item Exportar el modelo NLI a ONNX facilita unir investigación y
        producto: el backend sirve inferencia ligera y mantiene un
        fallback estable para CI y desarrollo local.
\end{itemize}
